\documentclass[UTF8,fontset=none,11pt,a4paper]{ctexart}
\usepackage{ipara-handbook}
\handbooksetup{NET-B03 Reliable Transfer × TCP}{408 · Network Internal Bridge}{Invariant · Byte Space · ACK · Timer}
\begin{document}
\handbooktitle{Reliable Transfer × TCP}{一般可靠性怎样落成 TCP 字节状态}{408 · NET-B03}

\begin{corebox}{Mother Interface}
\[
\boxed{\text{Delivery invariant}\to\text{TCP byte-state}\to\text{segment evidence}\to\text{event action}}
\]
NET03 给出“不重复、按序、最终恢复”的机制要求；NET06 把这些要求编码为每条连接两个方向的 SEQ/ACK、窗口、timer 与状态转换。
\end{corebox}

\begin{methodbox}{Owners 与停止边界}
NET03 Owns Seq/ACK/Timer/Window 的一般正确性与 GBN/SR 教材模型；NET06 Owns TCP byte stream、首部和连接状态。本册 Owns 两者的映射及类比停止条件。\code{rwnd/cwnd} 合成由 NET-B04，具体拥塞响应由 NET07。
\end{methodbox}

\section{从“帧编号”切换到“字节区间”}

教材 ARQ 常给每个 frame 一个离散序号；TCP SEQ 指向本 segment 第一个 data octet。若段携带 $L$ 字节且起始 SEQ 为 $s$，它占据半开区间 $[s,s+L)$；累计 ACK 为 $a$ 表示 $a$ 之前连续字节已收到、下一期待 $a$。SYN/FIN 各额外占一个序号，纯 ACK 通常不占。

因此映射必须以区间而非“第几个包”进行：
\[
\boxed{\text{Sent intervals}\supseteq\text{ACKed contiguous prefix}\quad\text{and}\quad
FlightSize=\text{sent but not cumulatively ACKed bytes}.}
\]

\section{观测不等于真相}

发送端不能直接看到“哪个 segment 在网络中丢了”，只能从 ACK、duplicate ACK、SACK（若题设给出）和 timer 推断。ACK 丢失可被后续累计 ACK 覆盖；数据重复到达必须被序号空间识别而不重复交付；乱序段是否缓存服从 TCP 状态/题设，不能把 GBN 的接收窗口 1 强套给 TCP。

\begin{center}\small
\begin{tabularx}{\linewidth}{L{2.4cm}Y Y}
\toprule
一般机制 & TCP 投影 & 边界\\
\midrule
identity & byte SEQ interval、connection tuple & 有限 32-bit 序号按模循环\\
positive evidence & cumulative ACK、可选 SACK & ACK 不表示应用已读取\\
loss suspicion & RTO、duplicate ACK pattern & 信号可能延迟/歧义\\
recovery action & retransmit missing bytes & 重传同时会触发 NET07 拥塞响应\\
duplicate suppression & receive-next 与已收区间 & 不得重复交付应用\\
\bottomrule
\end{tabularx}
\end{center}

\section{事件账本：一段丢失后的恢复}

发送方发 $[1000,1500)$、$[1500,2000)$、$[2000,2500)$。中段丢失而第三段先到，接收端连续前缀仍停在 1500，于是 ACK 不能跨过缺口。发送端收到重复 ACK 只是得到“1500 前存在缺口”的证据；达到题设 fast-retransmit 阈值或 RTO 到期后重传缺失区间。缺口补齐后，若第三段已缓存，累计 ACK 可一次推进到 2500。

这个例子解释三个常见现象：ACK 跳跃可确认多个已缓存区间；重传的是字节区间而不是“ACK number 对应的包”；RTO 与 duplicate ACK 是不同强度/时机的观测，具体拥塞窗口变化不由本桥决定。

\section{类比的可用范围}

GBN 帮助理解累计 ACK 和连续前缀阻塞；SR 帮助理解乱序缓存与选择性信息。但 TCP 不等于任一纯模型：它以 byte 为序号单位，可累计确认并缓存失序数据，还可能采用 SACK 和不同重传策略。类比一旦不能保持“同一状态、同一事件、同一不变量”，就应停止。

\section{调用协议与 First Divergence}

\begin{enumerate}
\item 写 connection tuple，并为两个方向分别建字节坐标；
\item 把每段换成 $[SEQ,SEQ+LEN+SYN+FIN)$；
\item 维护连续已确认前缀、乱序缓存与 in-flight 区间；
\item 事件到达时只更新它能证明的状态，不把推测写成事实；
\item 触发重传后，把拥塞窗口更新交给 NET07/NET-B04。
\end{enumerate}

\begin{warnbox}{Anti-Bridge}
TCP SEQ 每段加 1，是把 frame 序号硬套字节空间；ACK 到 $a$ 就说“$a$ 已收到”，是忘记 ACK 表示下一期望字节；收到重复 ACK 就断言网络拥塞，是把可靠性证据与拥塞判断合并；把 ACK 当应用消费证明，是越过传输交付边界。
\end{warnbox}

\begin{corebox}{Compression}
可靠性目标是什么？TCP 用哪个字节区间表示？当前 segment/ACK 能证明什么？哪个 timer 或证据触发动作？类比何时必须停止？
\end{corebox}
\end{document}
